% ===========================================================================
\chapter{Camada de Acesso a Dados e Persistência (Repository Layer)}
\label{cap_camada_persistencia}
% ===========================================================================

A camada de persistência é responsável por intermediar a comunicação entre o modelo orientado a objetos do SIGEA-GTT e o Sistema de Gerenciamento de Banco de Dados relacional PostgreSQL 16 LTS. Implementada com base no framework Spring Data JPA sobre o mecanismo do Hibernate ORM, esta camada segue rigorosamente o padrão *Repository* \cite{fowler2002}, abstraindo a geração de comandos SQL/DDL e provendo contratos de consulta fortemente tipados.

% ===========================================================================
\section{Diagrama de Classes da Camada de Repositórios}
\label{sec_diagrama_classes_repositorios}
% ===========================================================================

A Figura \ref{fig_diagrama_classes_repositorios} ilustra a hierarquia estrutural de interfaces que compõem a camada de acesso a dados, com ênfase na herança de \texttt{JpaRepository<T, UUID>} e na especificação das operações de busca especializadas:

\begin{figure}[!ht]
\centering
\caption{Diagrama de Classes de Repositórios (Spring Data JPA)}
\label{fig_diagrama_classes_repositorios}
\resizebox{0.96\textwidth}{!}{%
\begin{tikzpicture}[
    interface_box/.style={
        rectangle split,
        rectangle split parts=2,
        rectangle split part fill={teal!15, white},
        draw=teal!80!black,
        thick,
        rounded corners=2pt,
        font=\tiny,
        minimum width=5.2cm,
        text width=5.0cm,
        align=left
    },
    base_box/.style={
        rectangle split,
        rectangle split parts=2,
        rectangle split part fill={clinical-100!40, white},
        draw=clinical-700,
        thick,
        rounded corners=2pt,
        font=\tiny,
        minimum width=5.8cm,
        text width=5.6cm,
        align=left
    },
    inherit/.style={
        draw=clinical-900,
        thick,
        dashed,
        -{Triangle[open]}
    }
]

    % ====================================================
    % INTERFACE BASE JPA_REPOSITORY
    % ====================================================
    \node[base_box] (jpa_repo) at (3.5, 7.0) {
        \centering\textbf{\scriptsize <<Interface>>\\JpaRepository<T, ID>}
        \nodepart{two}
        + save(T) : T\\
        + findById(ID) : Optional<T>\\
        + findAll(Pageable) : Page<T>\\
        + delete(T) : void\\
        + count() : long
    };

    % ====================================================
    % USER REPOSITORY
    % ====================================================
    \node[interface_box] (user_repo) at (-3.5, 3.5) {
        \centering\textbf{\scriptsize <<Interface>>\\UserRepository}
        \nodepart{two}
        + findByEmail(String) : Optional<User>\\
        + existsByEmail(String) : boolean\\
        + findByRole(UserRole, Pageable) : Page<User>\\
        + findByIsActiveTrue(Pageable) : Page<User>
    };

    % ====================================================
    % ACADEMIC CLASS REPOSITORY
    % ====================================================
    \node[interface_box] (class_repo) at (3.5, 3.5) {
        \centering\textbf{\scriptsize <<Interface>>\\AcademicClassRepository}
        \nodepart{two}
        + findByProfessorId(UUID, Pageable) : Page<AcademicClass>\\
        + findByAcademicPeriod(String) : List<AcademicClass>\\
        + existsBySubjectNameAndClassCode(String, String, String) : boolean\\
        + findByStudentsId(UUID, Pageable) : Page<AcademicClass>
    };

    % ====================================================
    % ACTIVITY REPOSITORY
    % ====================================================
    \node[interface_box] (activity_repo) at (10.5, 3.5) {
        \centering\textbf{\scriptsize <<Interface>>\\ActivityRepository}
        \nodepart{two}
        + findByAcademicClassId(UUID) : List<Activity>\\
        + countByAcademicClassId(UUID) : long\\
        + findByIdAndAcademicClassId(UUID, UUID) : Optional<Activity>
    };

    % ====================================================
    % ACTIVITY SUBMISSION REPOSITORY
    % ====================================================
    \node[interface_box] (submission_repo) at (10.5, -1.0) {
        \centering\textbf{\scriptsize <<Interface>>\\ActivitySubmissionRepository}
        \nodepart{two}
        + findByActivityIdAndStudentId(UUID, UUID) : Optional<ActivitySubmission>\\
        + existsByActivityIdAndStudentId(UUID, UUID) : boolean\\
        + findByActivityId(UUID) : List<ActivitySubmission>\\
        + countByActivityAcademicClassIdAndGradeIsNull(UUID) : long\\
        + findByStudentId(UUID, Pageable) : Page<ActivitySubmission>
    };

    % ====================================================
    % GTT MODULE REPOSITORY
    % ====================================================
    \node[interface_box] (module_repo) at (-3.5, -1.0) {
        \centering\textbf{\scriptsize <<Interface>>\\GttModuleRepository}
        \nodepart{two}
        + findByCode(String) : Optional<GttModule>\\
        + existsByCode(String) : boolean\\
        + findAllByIsActiveTrue() : List<GttModule>
    };

    % ====================================================
    % GTT TRIGGER REPOSITORY
    % ====================================================
    \node[interface_box] (trigger_repo) at (3.5, -1.0) {
        \centering\textbf{\scriptsize <<Interface>>\\GttTriggerRepository}
        \nodepart{two}
        + findByCode(String) : Optional<GttTrigger>\\
        + findByModuleId(UUID) : List<GttTrigger>\\
        + existsByCode(String) : boolean\\
        + findAllByIsActiveTrue() : List<GttTrigger>
    };

    % ====================================================
    % HARM SEVERITY REPOSITORY
    % ====================================================
    \node[interface_box] (severity_repo) at (3.5, -4.8) {
        \centering\textbf{\scriptsize <<Interface>>\\HarmSeverityRepository}
        \nodepart{two}
        + findByCategoryLetter(String) : Optional<HarmSeverity>\\
        + findAllByOrderByCategoryLetterAsc() : List<HarmSeverity>\\
        + existsByCategoryLetter(String) : boolean
    };

    % ====================================================
    % RELAÇÕES DE HERANÇA (EXTENDS JPA_REPOSITORY)
    % ====================================================
    \draw[inherit] (user_repo.north) -- ++(0,0.8) -| (jpa_repo.south);
    \draw[inherit] (class_repo.north) -- (jpa_repo.south);
    \draw[inherit] (activity_repo.north) -- ++(0,0.8) -| (jpa_repo.south);
    \draw[inherit] (submission_repo.north) -- (activity_repo.south);
    \draw[inherit] (module_repo.north) -- (user_repo.south);
    \draw[inherit] (trigger_repo.north) -- (class_repo.south);
    \draw[inherit] (severity_repo.north) -- (trigger_repo.south);

\end{tikzpicture}
}
\fonte{Elaborado pelo autor (2026).}
\end{figure}

% ===========================================================================
\section{Detalhamento dos Contratos de Persistência}
\label{sec_detalhamento_repositorios}
% ===========================================================================

Cada repositório estende a interface genérica \texttt{JpaRepository<T, UUID>}, herdando rotinas de CRUD e paginação e definindo métodos de busca derivada (\textit{Query Methods}) baseados nas convenções do Spring Data:

\begin{enumerate}
    \item \textbf{\texttt{UserRepository}}:
    \begin{itemize}
        \item \texttt{findByEmail(String)}: Busca atômica por e-mail institucional, utilizada criticamente no fluxo de autenticação JWT;
        \item \texttt{existsByEmail(String)}: Verificação rápida de colisão cadastral antes de persistência;
        \item \texttt{findByRole(UserRole, Pageable)}: Listagem paginada segregando discentes, docentes e administradores com exatamente dez itens por página.
    \end{itemize}

    \item \textbf{\texttt{AcademicClassRepository}}:
    \begin{itemize}
        \item \texttt{findByProfessorId(UUID, Pageable)}: Filtra as turmas atribuídas a um professor titular específico;
        \item \texttt{existsBySubjectNameAndClassCodeAndAcademicPeriod(...)}: Valida a regra de unicidade composta para impedir duplicidade de turmas no mesmo período letivo semestral;
        \item \texttt{findByStudentsId(UUID, Pageable)}: Recupera as turmas nas quais determinado aluno possui matrícula ativa.
    \end{itemize}

    \item \textbf{\texttt{ActivitySubmissionRepository}}:
    \begin{itemize}
        \item \texttt{findByActivityIdAndStudentId(UUID, UUID)}: Localiza a submissão de um discente em determinada atividade, garantindo a verificação de tentativa única;
        \item \texttt{countByActivityAcademicClassIdAndGradeIsNull(}\linebreak \texttt{UUID)}: Consulta agregada que contabiliza submissões pendentes de nota na turma, utilizada como critério de bloqueio para o encerramento do semestre letivo (\texttt{UC10} / \texttt{RF13});
        \item \texttt{findByActivityId(UUID)}: Recupera todas as entregas para composição do painel de correção docente.
    \end{itemize}

    \item \textbf{\texttt{GttModuleRepository}, \texttt{GttTriggerRepository}, \texttt{HarmSeverityRepository}}:
    \begin{itemize}
        \item \texttt{findByCode(String)} / \texttt{findByCategoryLetter(String)}: Consultas exatas de alta velocidade aceleradas pelos índices B-Tree criados no banco de dados.
    \end{itemize}
\end{enumerate}

% ===========================================================================
\section{Políticas Transacionais e Concorrência}
\label{sec_politicas_transacionais}
% ===========================================================================

A camada de repositório opera em estreita consonância com o gerenciador de transações do Spring Framework (\texttt{PlatformTransactionManager}). No SIGEA-GTT, adotam-se as seguintes diretrizes transacionais:
\begin{itemize}
    \item \textbf{Transações Somente-Leitura (\texttt{@Transactional(readOnly = true)})}: Aplicadas em todos os métodos de consulta para desabilitar o mecanismo de \textit{dirty checking} do Hibernate, reduzindo o consumo de memória da JVM e acelerando a resposta HTTP;
    \item \textbf{Transações Mutativas Atômicas} (\texttt{@Transactional} com \texttt{READ\_COMMITTED}): Utilizadas em criação de turmas, matrículas e submissões clínicas para assegurar as propriedades ACID e integridade referencial frente a acessos concorrentes na rede.
\end{itemize}
